% =================================================================
% Packages and macros for the Wiley Systems Engineering article (syseng.cls).
% Document class, authors, and abstract live in main.tex.
% =================================================================

\graphicspath{{./}{figures/}}
\usepackage[space]{grffile}
\usepackage{textcomp}
\usepackage{longtable}
\usepackage{colortbl}
\usepackage{tabularx}
\usepackage{multirow}
\usepackage{mathtools}
\usepackage{bm}
\usepackage{etoolbox}
\usepackage[english]{babel}
\usepackage{siunitx}
\usepackage{placeins}

\usepackage{tikz}
\usetikzlibrary{arrows,arrows.meta,positioning,shapes,shapes.geometric,fit,backgrounds,calc,automata}

% Compatibility aliases for natbib-style citation commands used in the text.
\providecommand{\citep}[1]{\cite{#1}}
\providecommand{\citet}[1]{\cite{#1}}
\providecommand{\citealt}[1]{\cite{#1}}

% Theorem and proof environments matching Wiley Systems Engineering style.
\newtheorem{theorem}{Theorem}
\newtheorem{lemma}[theorem]{Lemma}
\newtheorem{proposition}[theorem]{Proposition}
\newtheorem{corollary}[theorem]{Corollary}

\makeatletter
\def\@opargbegintheorem#1#2#3{\trivlist
   \item[\hskip \labelsep{\bfseries #1\ #2}]\textbf{(#3).}\ \itshape}
\makeatother

\newenvironment{proof}[1][Proof]{\begin{trivlist}
\item[\hskip \labelsep {\bfseries #1.}]}{\end{trivlist}}
\newenvironment{definition}[1][Definition]{\begin{trivlist}
\item[\hskip \labelsep {\bfseries #1.}]}{\end{trivlist}}
\newenvironment{example}[1][Example]{\begin{trivlist}
\item[\hskip \labelsep {\bfseries #1.}]}{\end{trivlist}}
\newenvironment{remark}[1][Remark]{\begin{trivlist}
\item[\hskip \labelsep {\bfseries #1.}]}{\end{trivlist}}

\newcommand{\qed}{\nobreak \ifvmode \relax \else
      \ifdim\lastskip<1.5em \hskip-\lastskip
      \hskip1.5em plus0em minus0.5em \fi \nobreak
      \vrule height0.75em width0.5em depth0.25em\fi}

% Table styling helpers.
\definecolor{tableheader}{HTML}{E6E6E6}
\newcolumntype{Y}{>{\raggedright\arraybackslash}X}
\newcolumntype{P}[1]{>{\centering\arraybackslash}p{#1}}
\newcolumntype{A}[1]{>{\raggedright\arraybackslash}p{#1}}
\newcolumntype{M}[1]{>{\raggedright\arraybackslash}m{#1}}
\newcommand{\headingfont}{}

% Lightweight stand-in for mini-heading macro.
\newcommand{\miniheading}[1]{%
  \par\vspace{4bp}\noindent\textbf{#1.}\space
}

% Safe figure include (optional assets).
\newcommand{\colfig}[2][]{%
  \IfFileExists{#2}{\includegraphics[width=\linewidth,#1]{#2}}{%
    \fbox{\parbox[b][1.5in][c]{\linewidth}{\centering \textit{Missing figure: }#2}}}%
}

\newif\iflatexml\latexmlfalse
\providecommand{\tightlist}{\setlength{\itemsep}{0pt}\setlength{\parskip}{0pt}}%

\AtBeginDocument{\DeclareGraphicsExtensions{.pdf,.PDF,.eps,.EPS,.png,.PNG,.tif,.TIF,.jpg,.JPG,.jpeg,.JPEG}}
